\documentclass[a4paper,10pt,brazil]{article}
\usepackage{provastyle}
\usepackage{menukeys}

\begin{document}
\makeexamtitlepage{UFOP}{BCC222: Programação Funcional}{José Romildo}{2026/1}{Primeira Prova}{01}
\begin{multicols}{2}
\questionTitle{1}{read isolado no GHCi com EDR}
\begin{flushright}
  \footnotesize
  \textit{\textbf{7: Polimorfismo}}\\
  \textit{cap07-defaulting-read-isolado}\\

\end{flushright}
No GHCi com \haskellinline{ExtendedDefaultRules} ligada, o que acontece conceitualmente com a expressão abaixo?

\begin{minted}{haskell}
read "23"
\end{minted}

\begin{enumerate}
  \item O tipo é inferido como \haskellinline{Int} pelo conteúdo textual \haskellinline{"23"}.
  \item O tipo é inferido como \haskellinline{Integer} pela regra padrão do Relatório, sem depender do GHCi.
  \item O GHCi adiciona uma restrição de impressão, tenta defaultar o tipo para \haskellinline{()}, e a análise da string \haskellinline{"23"} falha em tempo de execução.
  \item A expressão é aceita e retorna a string original, sem parse.
  \item A expressão tem tipo \haskellinline{String}, pois \haskellinline{read} nunca converte valores.

\end{enumerate}

\questionTitle{2}{A onipresença do Prelude}
\begin{flushright}
  \footnotesize
  \textit{\textbf{2: Ambiente de desenvolvimento}}\\
  \textit{c02-ghci-modulos-01}\\

\end{flushright}
No ambiente GHCi, é possível utilizar funções matemáticas como \haskellinline{sqrt} ou operadores como \haskellinline{+} imediatamente após abrir o terminal, sem a necessidade de comandos adicionais. Qual é a justificativa técnica para isso?

\begin{enumerate}
  \item O GHCi possui um recurso de compilação JIT que adivinha a intenção do programador e cria as funções matemáticas em tempo de execução.

  \item O GHCi carrega automaticamente, por padrão, o módulo especial \texttt{Prelude}, que contém um grande conjunto de funções e tipos básicos.

  \item Linguagens funcionais não requerem bibliotecas de base, pois todas as operações aritméticas são resolvidas diretamente pelo \emph{hardware} da CPU.

  \item O comando \texttt{ghci} faz uma varredura no disco rígido buscando qualquer arquivo \texttt{.hs} pré-existente e importa todas as funções encontradas silenciosamente.

  \item Essas funções não pertencem a nenhum módulo; elas são palavras-chave imutáveis encravadas no analisador léxico (\emph{parser}) do compilador.


\end{enumerate}

\questionTitle{3}{Recursão e funções de redução}
\begin{flushright}
  \footnotesize
  \textit{\textbf{9: Funções recursivas}}\\
  \textit{c09-eficiencia-reducao-01}\\

\end{flushright}
O capítulo observa que, em linguagens funcionais, repetições podem ser expressas por recursão ou por abstrações construídas sobre padrões recursivos. Qual alternativa expressa melhor essa ideia?

\begin{enumerate}
  \item Funções de redução são equivalentes a variáveis globais atualizadas em laços.

  \item Funções de redução substituem completamente a necessidade de compreender casos base.

  \item Funções de redução, como \haskellinline|foldr| e \haskellinline|foldl|, encapsulam padrões recorrentes de recursão sobre estruturas.

  \item Funções de redução só existem em linguagens imperativas.

  \item Funções de redução impedem o uso de funções puras.


\end{enumerate}

\questionTitle{4}{O ciclo prático de edição e recarregamento}
\begin{flushright}
  \footnotesize
  \textit{\textbf{2: Ambiente de desenvolvimento}}\\
  \textit{c02-fluxo-trabalho-01}\\

\end{flushright}
Qual das seguintes sequências descreve adequadamente o fluxo de trabalho (\emph{workflow}) mais ágil e recomendado para o desenvolvimento de programas locais usando Haskell e GHCi?

\begin{enumerate}
  \item Editar o código-fonte em um editor externo e salvá-lo $\rightarrow$ Usar o comando \texttt{:reload} (ou \texttt{:r}) no GHCi já aberto $\rightarrow$ Testar a nova função iterativamente.

  \item Digitar o comando \texttt{:load} no editor VS Code $\rightarrow$ Modificar a função $\rightarrow$ Enviar a compilação final para o GitHub Classroom rodar as ferramentas de análise.

  \item Editar as funções diretamente dentro do GHCi através da declaração múltipla $\rightarrow$ Encerrar o GHCi com \texttt{:quit} $\rightarrow$ Reabrir o terminal para que as alterações sejam gravadas no disco.

  \item Reescrever todo o código no GHCi $\rightarrow$ Digitar o comando de exportação de binário $\rightarrow$ Rodar o \emph{script} nativo em C.

  \item Modificar o arquivo \texttt{.hs} no editor $\rightarrow$ Fechar o GHCi $\rightarrow$ Invocar o compilador GHC na linha de comando para gerar um novo pacote executável $\rightarrow$ Testar o binário.


\end{enumerate}

\questionTitle{5}{Significado do operador ::}
\begin{flushright}
  \footnotesize
  \textit{\textbf{4: Tipos de Dados}}\\
  \textit{c04-assinatura-operador-01}\\

\end{flushright}
Na declaração \haskellinline{minhaFuncao :: Int -> Bool}, o que o operador \haskellinline{::} significa?

\begin{enumerate}
  \item \emph{pertence à classe}

  \item \emph{tem o tipo de}

  \item \emph{é definido como}

  \item \emph{retorna}

  \item \emph{é um subtipo de}


\end{enumerate}

\questionTitle{6}{Avaliação direta de if}
\begin{flushright}
  \footnotesize
  \textit{\textbf{5: Expressão condicional}}\\
  \textit{c05-if-avaliacao-01}\\

\end{flushright}
Qual é o valor da expressão abaixo?
\begin{haskellcode}
if null [] then 'a' else 'b'
\end{haskellcode}

\begin{enumerate}
  \item \haskellinline|'b'|

  \item \haskellinline|'a'|

  \item \haskellinline|"a"|

  \item \haskellinline|True|

  \item Ocorre um erro de sintaxe.


\end{enumerate}

\questionTitle{7}{Definição formal de transparência referencial}
\begin{flushright}
  \footnotesize
  \textit{\textbf{1: Paradigmas de programação}}\\
  \textit{c01-transparencia-referencial-01}\\

\end{flushright}
Na teoria de linguagens de programação funcionais, o que significa afirmar que uma dada expressão possui \textbf{transparência referencial}?

\begin{enumerate}
  \item Significa que o sistema de inferência de tipos consegue deduzir a assinatura da função mesmo que o programador omita essa informação no cabeçalho.

  \item Significa que a expressão pode ser substituída diretamente pelo seu valor computado correspondente em qualquer ponto do programa, sem alterar o comportamento global do sistema.

  \item Significa que a função tem permissão explícita para ler variáveis do escopo global mutável, desde que não as modifique durante a execução.

  \item Significa que a função permite que referências de ponteiros sejam passadas como argumentos em vez de cópias de valores, economizando alocação de memória.

  \item Significa que o código-fonte da função não pode ser ofuscado pelo compilador, permanecendo \emph{transparente} e legível para fins de auditoria no binário.


\end{enumerate}

\questionTitle{8}{Inferência com zip, head e fst}
\begin{flushright}
  \footnotesize
  \textit{\textbf{7: Polimorfismo}}\\
  \textit{cap07-param-zip-fst-inferencia}\\

\end{flushright}
Considere:

\begin{minted}{haskell}
f xs ys = fst (head (zip xs ys))
\end{minted}

Qual é a assinatura de tipo mais geral de \haskellinline{f}?

\begin{enumerate}
  \item \mintinline{haskell}{f :: Ord a => [a] -> [b] -> a}
  \item \mintinline{haskell}{f :: [a] -> [a] -> a}
  \item \mintinline{haskell}{f :: [a] -> [b] -> b}
  \item \mintinline{haskell}{f :: [(a,b)] -> a}
  \item \mintinline{haskell}{f :: [a] -> [b] -> a}

\end{enumerate}

\questionTitle{9}{Erro por desalinhamento na regra de layout}
\begin{flushright}
  \footnotesize
  \textit{\textbf{3: Expressões e definições}}\\
  \textit{c03-layout-avaliacao-01}\\

\end{flushright}
Considere o trecho abaixo, no qual a segunda definição local não está alinhada com a primeira.

\begin{haskellcode}
f x = let a = x + 1
        b = x + 2
      in a + b
\end{haskellcode}

Qual é a causa do problema de sintaxe?

\begin{enumerate}
  \item Pela regra de layout, definições pertencentes ao mesmo bloco devem começar na mesma coluna; \haskellinline|b| deveria estar alinhado com \haskellinline|a|.

  \item O problema é que \haskellinline|in| deve aparecer antes de todas as definições do bloco.

  \item O problema é o uso de \haskellinline|+| dentro de definições locais.

  \item O problema é que Haskell ignora espaços em branco, então o compilador não consegue separar tokens.

  \item O problema é o uso de \haskellinline|let|; definições locais só podem ser feitas com \haskellinline|where|.


\end{enumerate}

\questionTitle{10}{Cláusula else obrigatória}
\begin{flushright}
  \footnotesize
  \textit{\textbf{5: Expressão condicional}}\\
  \textit{c05-if-regras-01}\\

\end{flushright}
Em Haskell, por que a expressão abaixo é incorreta?
\begin{haskellcode}
if x > 0 then x
\end{haskellcode}

\begin{enumerate}
  \item Porque a condição \haskellinline|x > 0| deveria estar entre parênteses.

  \item Porque o ramo \haskellinline|then| deveria conter uma expressão mais completa do que \haskellinline|x|.

  \item Porque falta o ramo \haskellinline|else|, obrigatório para que a expressão sempre produza um valor.

  \item Porque o ramo \haskellinline|then| só pode conter valores booleanos.

  \item Porque expressões condicionais só podem ser usadas dentro do corpo de funções nomeadas.


\end{enumerate}

\questionTitle{11}{Sintaxe de aplicação prefixa}
\begin{flushright}
  \footnotesize
  \textit{\textbf{3: Expressões e definições}}\\
  \textit{c03-aplicacao-notacoes-01}\\

\end{flushright}
Considere a expressão \haskellinline|mod 25 7|. Qual interpretação está correta de acordo com a sintaxe de aplicação de funções em Haskell?

\begin{enumerate}
  \item A expressão aplica \haskellinline|25| como função aos argumentos \haskellinline|mod| e \haskellinline|7|.

  \item A expressão é inválida porque funções em Haskell sempre exigem parênteses em volta dos argumentos.

  \item A função \haskellinline|mod| é aplicada a dois argumentos, \haskellinline|25| e \haskellinline|7|, separados por espaços.

  \item A expressão só seria válida se fosse escrita como \haskellinline|mod(25, 7)|.

  \item \haskellinline|mod| recebe um único argumento, a tupla \haskellinline|(25,7)|.


\end{enumerate}

\questionTitle{12}{Finalidade da cláusula where}
\begin{flushright}
  \footnotesize
  \textit{\textbf{3: Expressões e definições}}\\
  \textit{c03-definicoes-where-01}\\

\end{flushright}
Qual alternativa descreve corretamente o papel da cláusula \haskellinline|where| em uma definição Haskell?

\begin{enumerate}
  \item Ela substitui a necessidade do sinal de igualdade em definições de funções.

  \item Ela introduz uma expressão que pode ser escrita livremente em qualquer posição onde um literal seria permitido.

  \item Ela anexa definições locais a uma equação, permitindo nomear subexpressões auxiliares no escopo dessa definição.

  \item Ela torna todas as definições auxiliares visíveis globalmente no módulo.

  \item Ela serve exclusivamente para importar nomes de outros módulos.


\end{enumerate}

\questionTitle{13}{Comando :type no GHCi}
\begin{flushright}
  \footnotesize
  \textit{\textbf{4: Tipos de Dados}}\\
  \textit{c04-ghci-consulta-01}\\

\end{flushright}
Qual comando do GHCi exibe o tipo de uma expressão sem avaliá-la?

\begin{enumerate}
  \item \texttt{:info}

  \item \texttt{:show}

  \item \texttt{:kind}

  \item \texttt{:type}

  \item \texttt{:load}


\end{enumerate}

\questionTitle{14}{getChar e getLine}
\begin{flushright}
  \footnotesize
  \textit{\textbf{8: Programas interativos}}\\
  \textit{c08-entrada-saida-05}\\

\end{flushright}
Qual alternativa compara corretamente \haskellinline|getChar| e \haskellinline|getLine|?

\begin{enumerate}
  \item \haskellinline|getChar| é pura; \haskellinline|getLine| é uma ação de E/S.

  \item \haskellinline|getLine| recebe uma string de prompt como argumento, enquanto \haskellinline|getChar| não recebe argumentos.

  \item \haskellinline|getChar| imprime um caractere; \haskellinline|getLine| imprime uma linha.

  \item \haskellinline|getChar| tem tipo \haskellinline|IO Char| e lê um caractere; \haskellinline|getLine| tem tipo \haskellinline|IO String| e lê uma linha de texto.

  \item Ambas têm tipo \haskellinline|IO String|, mas \haskellinline|getChar| retorna uma string de tamanho 1.


\end{enumerate}

\questionTitle{15}{Separação entre leitura e soma pura}
\begin{flushright}
  \footnotesize
  \textit{\textbf{10: Ações de E/S recursivas}}\\
  \textit{c10-separacao-puro-01}\\

\end{flushright}
Considere a versão que separa a leitura da soma:

\begin{haskellcode}
main :: IO ()
main = do
  lista <- lerLista
  print (sum lista)

lerLista :: IO [Integer]
lerLista = do
  x <- readLn
  if x == 0
    then return []
    else do
      resto <- lerLista
      return (x:resto)
\end{haskellcode}

Qual é a principal vantagem dessa organização?

\begin{enumerate}
  \item Ela faz com que \haskellinline|sum| leia diretamente os números do teclado.

  \item Ela evita a necessidade de um caso base, pois \haskellinline|sum| detecta automaticamente o sentinela.

  \item Ela elimina completamente o uso de \haskellinline|IO| do programa.

  \item Ela obriga o compilador a transformar \haskellinline|lerLista| em um laço imperativo.

  \item Ela isola a interação com o usuário em \haskellinline|lerLista| e permite que o processamento, como \haskellinline|sum lista|, permaneça puro.


\end{enumerate}

\questionTitle{16}{Comprimentos diferentes na correção de provas}
\begin{flushright}
  \footnotesize
  \textit{\textbf{10: Ações de E/S recursivas}}\\
  \textit{c10-bordas-robustez-05}\\

\end{flushright}
Considere uma função pura auxiliar para corrigir uma prova de múltipla escolha:

\begin{haskellcode}
corrige :: String -> String -> Int
corrige [] [] = 0
corrige (g:gs) (r:rs)
  | g == r    = 1 + corrige gs rs
  | otherwise = corrige gs rs
\end{haskellcode}

Qual problema pode ocorrer se o gabarito e as respostas tiverem comprimentos diferentes?

\begin{enumerate}
  \item A função passa a ter tipo \haskellinline|IO Int|, pois listas de tamanhos diferentes exigem entrada e saída.

  \item A função sempre considera respostas ausentes como corretas, pois \haskellinline|[]| vale \haskellinline|True|.

  \item O compilador rejeita a definição antes da execução, pois não é permitido casar padrões com duas listas.

  \item A função pode falhar por padrões não exaustivos, pois não há casos para \haskellinline|[]| contra lista não vazia ou lista não vazia contra \haskellinline|[]|.

  \item A função inverte automaticamente a maior lista para tentar alinhar os elementos.


\end{enumerate}

\questionTitle{17}{Visibilidade em guardas e resultados}
\begin{flushright}
  \footnotesize
  \textit{\textbf{5: Expressão condicional}}\\
  \textit{c05-where-escopo-01}\\

\end{flushright}
Em uma equação com guardas e uma cláusula \haskellinline|where|, onde são visíveis as definições locais?

\begin{enumerate}
  \item Apenas na primeira guarda da equação.

  \item Apenas nas guardas, nunca nas expressões de resultado.

  \item Em toda a função, inclusive em outras equações com o mesmo nome.

  \item Apenas nas expressões de resultado, nunca nas guardas.

  \item Nas guardas e nas expressões de resultado da mesma equação.


\end{enumerate}

\questionTitle{18}{Avaliação preguiçosa e acumuladores}
\begin{flushright}
  \footnotesize
  \textit{\textbf{10: Ações de E/S recursivas}}\\
  \textit{c10-acumuladores-cauda-05}\\

\end{flushright}
No capítulo, a versão com acumulador é apresentada como uma melhoria estrutural em relação à versão que deixa somas pendentes. Em Haskell, qual ressalva técnica é importante ao discutir uso de memória nessa versão?

\begin{enumerate}
  \item A ressalva é que \haskellinline|readLn| avalia todos os valores futuros antes de retornar o primeiro.

  \item A ressalva é que \haskellinline|return| sempre força completamente o valor acumulado antes de empacotá-lo em \haskellinline|IO|.

  \item Mesmo com chamada em posição de cauda, a avaliação preguiçosa pode acumular expressões não avaliadas no acumulador se ele não for forçado adequadamente.

  \item A ressalva é que toda função recursiva de cauda em Haskell é rejeitada pelo compilador por não ser declarativa.

  \item A ressalva é que acumuladores não podem ser usados em ações de E/S, apenas em funções puras.


\end{enumerate}

\questionTitle{19}{Consulta com zip e lookup}
\begin{flushright}
  \footnotesize
  \textit{\textbf{6: Estruturas de dados básicas}}\\
  \textit{c06-lookup-associacao-01}\\

\end{flushright}
Considere a expressão abaixo.
\begin{minted}{haskell}
let chaves = [2,4,6,8]
    valores = ["m","n","o","p"]
    tabela = zip chaves valores
in lookup 8 tabela
\end{minted}
Qual é o resultado final?

\begin{enumerate}
  \item \haskellinline{ "p" }

  \item \haskellinline{ Nothing }

  \item \haskellinline{ Just "o" }

  \item \haskellinline{ Just "p" }

  \item Ocorre um erro de tipo.


\end{enumerate}

\questionTitle{20}{Inferência de tipo de not True}
\begin{flushright}
  \footnotesize
  \textit{\textbf{4: Tipos de Dados}}\\
  \textit{c04-inferencia-basica-01}\\

\end{flushright}
Qual é o tipo inferido para a expressão \haskellinline{not True}?

\begin{enumerate}
  \item \haskellinline{Boolean}

  \item \haskellinline{Bool}

  \item \haskellinline{False}

  \item \haskellinline{True}

  \item \haskellinline{Bool -> Bool}


\end{enumerate}

\questionTitle{21}{Declaração de sinônimo de tipo}
\begin{flushright}
  \footnotesize
  \textit{\textbf{4: Tipos de Dados}}\\
  \textit{c04-sinonimos-criacao-01}\\

\end{flushright}
Qual palavra-chave é usada para criar um sinônimo de tipo em Haskell?

\begin{enumerate}
  \item \haskellinline{newtype}

  \item \haskellinline{class}

  \item \haskellinline{type}

  \item \haskellinline{instance}

  \item \haskellinline{data}


\end{enumerate}

\questionTitle{22}{Tupla, lista e agrupamento}
\begin{flushright}
  \footnotesize
  \textit{\textbf{6: Estruturas de dados básicas}}\\
  \textit{c06-tuplas-sintaxe-01}\\

\end{flushright}
Qual das expressões abaixo representa uma tupla de exatamente dois elementos?

\begin{enumerate}
  \item \haskellinline{(42)}

  \item \haskellinline{[42, "ok"]}

  \item \haskellinline{(42, "ok")}

  \item \haskellinline{("ok")}

  \item \haskellinline{42, "ok"}


\end{enumerate}

\questionTitle{23}{Necessidade de reverse ao acumular com cons}
\begin{flushright}
  \footnotesize
  \textit{\textbf{10: Ações de E/S recursivas}}\\
  \textit{c10-listas-ordem-01}\\

\end{flushright}
Considere a função:

\begin{haskellcode}
lerLista :: [Integer] -> IO [Integer]
lerLista xs = do
  x <- readLn
  if x == 0
    then return (reverse xs)
    else lerLista (x:xs)
\end{haskellcode}

Por que \haskellinline|reverse| é usado no caso base?

\begin{enumerate}
  \item Porque \haskellinline|reverse| força a ação de E/S a ser executada antes do caso base.

  \item Porque cada novo valor é inserido no início do acumulador com \haskellinline|x:xs|, invertendo a ordem de entrada.

  \item Porque \haskellinline|sum| exige que a lista esteja invertida para funcionar corretamente.

  \item Porque \haskellinline|readLn| sempre lê os números na ordem inversa em que foram digitados.

  \item Porque \haskellinline|return| só aceita listas ordenadas em ordem crescente.


\end{enumerate}

\questionTitle{24}{Forma fundamental de uma lista}
\begin{flushright}
  \footnotesize
  \textit{\textbf{6: Estruturas de dados básicas}}\\
  \textit{c06-listas-estrutura-01}\\

\end{flushright}
Qual é a forma fundamental da lista \haskellinline{['a', 'b', 'c']} usando apenas o construtor \haskellinline{(:)} e a lista vazia?

\begin{enumerate}
  \item \haskellinline{[] : 'a' : 'b' : 'c'}

  \item \haskellinline{('a', 'b', 'c')}

  \item \haskellinline{['a'] : ['b'] : ['c']}

  \item \haskellinline{'a' : 'b' : 'c'}

  \item \haskellinline{'a' : ('b' : ('c' : []))}


\end{enumerate}

\questionTitle{25}{Valor puro e ação de E/S}
\begin{flushright}
  \footnotesize
  \textit{\textbf{8: Programas interativos}}\\
  \textit{c08-tipos-acoes-05}\\

\end{flushright}
Qual alternativa descreve corretamente a diferença entre \haskellinline|"abc"| e \haskellinline|getLine|?

\begin{enumerate}
  \item \haskellinline|"abc"| tem tipo \haskellinline|IO ()| e \haskellinline|getLine| tem tipo \haskellinline|String|.

  \item Ambas têm tipo \haskellinline|IO String|, pois toda string em Haskell pode ser impressa no terminal.

  \item \haskellinline|getLine| é uma função pura do tipo \haskellinline|() -> String|, equivalente a uma função sem argumentos.

  \item Ambas têm tipo \haskellinline|String|, mas \haskellinline|getLine| é avaliada de forma preguiçosa.

  \item \haskellinline|"abc"| é uma \haskellinline|String| pura; \haskellinline|getLine| é uma ação do tipo \haskellinline|IO String| que pode produzir uma \haskellinline|String| quando executada.


\end{enumerate}

\questionTitle{26}{Restrição necessária para sort}
\begin{flushright}
  \footnotesize
  \textit{\textbf{7: Polimorfismo}}\\
  \textit{cap07-expressao-sort-restricao}\\

\end{flushright}
A função \haskellinline{sort} tem, conceitualmente, tipo:

\begin{minted}{haskell}
sort :: Ord a => [a] -> [a]
\end{minted}

Por que a restrição \haskellinline{Ord a} aparece?

\begin{enumerate}
  \item Porque qualquer função que recebe uma lista precisa de \haskellinline{Ord}.
  \item Porque \haskellinline{sort} converte todos os elementos para \haskellinline{String}.
  \item Porque \haskellinline{Ord} é necessário para construir listas, mesmo sem ordenação.
  \item Porque ordenar exige comparar elementos, e essas comparações dependem de uma instância de \haskellinline{Ord} para o tipo dos elementos.
  \item Porque \haskellinline{sort} altera os elementos da lista usando aritmética.

\end{enumerate}

\questionTitle{27}{Fundação teórica: Cálculo Lambda}
\begin{flushright}
  \footnotesize
  \textit{\textbf{1: Paradigmas de programação}}\\
  \textit{c01-historia-evolucao-01}\\

\end{flushright}
Na década de 1930, antes mesmo da existência dos computadores eletrônicos construídos em silício, foi desenvolvido um sistema formal matemático que provou ser tão poderoso e universal quanto a Máquina de Turing. Este sistema lançou a fundação teórica de toda a programação funcional moderna.

Qual é o nome deste sistema matemático e quem foi o seu criador?

\begin{enumerate}
  \item A Arquitetura Pura, desenvolvida por John von Neumann.

  \item O Cálculo Lambda, desenvolvido por Alonzo Church.

  \item A Álgebra Booleana, desenvolvida por George Boole.

  \item A Teoria da Avaliação Preguiçosa (\emph{Lazy Evaluation}), formalizada por David Turner.

  \item O Sistema de Inferência de Tipos de Primeira Ordem, desenvolvido por Robin Milner.


\end{enumerate}

\questionTitle{28}{Significado de otherwise}
\begin{flushright}
  \footnotesize
  \textit{\textbf{5: Expressão condicional}}\\
  \textit{c05-otherwise-01}\\

\end{flushright}
O que \haskellinline|otherwise| representa em uma definição com guardas?

\begin{enumerate}
  \item Um nome predefinido equivalente a \haskellinline|True|, normalmente usado como caso padrão.

  \item Uma variável criada automaticamente com o valor da guarda anterior.

  \item Um nome reservado equivalente a \haskellinline|False|.

  \item Uma forma especial de indicar erro em tempo de execução.

  \item Uma palavra reservada que faz a avaliação voltar para a equação anterior.


\end{enumerate}

\questionTitle{29}{Quantidade de termos na aproximação de pi}
\begin{flushright}
  \footnotesize
  \textit{\textbf{10: Ações de E/S recursivas}}\\
  \textit{c10-rastreamento-05}\\

\end{flushright}
Considere a versão revisada da aproximação de \(\pi\):

\begin{haskellcode}
valorPi :: Integer -> Double
valorPi n
  | n <= 0    = 0
  | otherwise = 4 * soma 0 (n - 1)

soma :: Integer -> Integer -> Double
soma i n
  | i <= n    = (-1) ^ i / (2 * fromIntegral i + 1)
                  + soma (i + 1) n
  | otherwise = 0
\end{haskellcode}

Se \haskellinline|valorPi 4| é avaliada, quais valores de \haskellinline|i| são usados nos termos da série?

\begin{enumerate}
  \item Nenhum valor, pois \haskellinline|valorPi 4| cai no caso \haskellinline|n <= 0|.

  \item \haskellinline|0, 1, 2, 3, 4|, pois a função sempre usa \haskellinline|n + 1| termos.

  \item \haskellinline|1, 2, 3, 4|, pois o parâmetro indica o maior índice.

  \item \haskellinline|4| apenas, pois \haskellinline|soma| avalia somente o último termo.

  \item \haskellinline|0, 1, 2, 3|, pois o parâmetro indica a quantidade de termos e o limite superior é \haskellinline|n - 1|.


\end{enumerate}

\questionTitle{30}{O comando :type e sua utilidade}
\begin{flushright}
  \footnotesize
  \textit{\textbf{2: Ambiente de desenvolvimento}}\\
  \textit{c02-ghci-inspecao-01}\\

\end{flushright}
Durante uma sessão interativa no GHCi, o programador utiliza frequentemente o comando \texttt{:type} (ou \texttt{:t}). O que este comando executa quando aplicado a uma expressão?

\begin{enumerate}
  \item Ele exibe o código-fonte original (em linguagem Haskell) que originou a função solicitada.

  \item Ele inspeciona e exibe a assinatura de tipo inferida pelo compilador para a expressão fornecida, sem efetivamente avaliá-la (executá-la).

  \item Ele pesquisa no repositório oficial da linguagem quais pacotes exportam tipos de dados que correspondam à expressão digitada.

  \item Ele força a avaliação completa da expressão matemática e formata a saída em um texto de múltiplas linhas legível para o usuário.

  \item Ele modifica o tipo da variável em tempo real, permitindo que um inteiro seja tratado como \emph{string} no GHCi temporariamente.


\end{enumerate}

\questionTitle{31}{Quando usar realToFrac}
\begin{flushright}
  \footnotesize
  \textit{\textbf{7: Polimorfismo}}\\
  \textit{cap07-conversao-realToFrac}\\

\end{flushright}
Qual é o papel de \haskellinline{realToFrac} em conversões numéricas?
\begin{enumerate}
  \item Substituir \haskellinline{read} na conversão de \haskellinline{String} para número.
  \item Permitir que operadores como \haskellinline{(+)} misturem automaticamente \haskellinline{Int} e \haskellinline{Double}.
  \item Arredondar qualquer número fracionário para o inteiro mais próximo.
  \item Converter de um tipo instância de \haskellinline{Real} para um tipo de destino instância de \haskellinline{Fractional}, como de \haskellinline{Float} para \haskellinline{Double}.
  \item Converter apenas de \haskellinline{Int} para \haskellinline{Integer}, preservando integralidade.

\end{enumerate}

\questionTitle{32}{Rastreamento de Fibonacci}
\begin{flushright}
  \footnotesize
  \textit{\textbf{9: Funções recursivas}}\\
  \textit{c09-formas-fibonacci-rastreamento-01}\\

\end{flushright}
Usando a definição do capítulo,
\begin{haskellcode}
fib n
  | n == 0 = 0
  | n == 1 = 1
  | n > 1  = fib (n - 2) + fib (n - 1)
\end{haskellcode}
qual é o valor de \haskellinline|fib 4|?

\begin{enumerate}
  \item \haskellinline|8|

  \item \haskellinline|5|

  \item \haskellinline|3|

  \item \haskellinline|2|

  \item \haskellinline|4|


\end{enumerate}

\questionTitle{33}{A ideia de mundo em IO}
\begin{flushright}
  \footnotesize
  \textit{\textbf{8: Programas interativos}}\\
  \textit{c08-modelo-io-05}\\

\end{flushright}
No modelo conceitual apresentado para programas interativos em Haskell, qual é a melhor interpretação da ideia de \emph{mundo}?

\begin{enumerate}
  \item É uma estrutura de dados concreta do Prelude que o programador deve passar manualmente para cada função.

  \item É apenas outro nome para o tipo \haskellinline|String| usado nas entradas do usuário.

  \item É uma forma conceitual de representar o ambiente externo com o qual o programa interage, como terminal, arquivos e outros dispositivos, preservando a separação entre valores puros e ações de E/S.

  \item É um módulo obrigatório que precisa ser importado para usar \haskellinline|getLine| e \haskellinline|putStrLn|.

  \item É uma variável global chamada \haskellinline|World| que pode ser lida e modificada por qualquer expressão pura.


\end{enumerate}

\questionTitle{34}{Sintaxe correta com guardas}
\begin{flushright}
  \footnotesize
  \textit{\textbf{5: Expressão condicional}}\\
  \textit{c05-guardas-sintaxe-semantica-01}\\

\end{flushright}
Qual das definições abaixo está escrita corretamente usando guardas em Haskell?

\begin{enumerate}
  \item \begin{haskellcode}
abs' n
  n < 0 -> -n
  otherwise -> n
\end{haskellcode}

  \item \begin{haskellcode}
abs' n
  if n < 0 = -n
  else     = n
\end{haskellcode}

  \item \begin{haskellcode}
abs' n
  | n < 0     = -n
  | otherwise = n
\end{haskellcode}

  \item \begin{haskellcode}
abs' n
  | n < 0 then -n
  | otherwise then n
\end{haskellcode}

  \item \begin{haskellcode}
abs' n =
  | n < 0     = -n
  | otherwise = n
\end{haskellcode}


\end{enumerate}

\questionTitle{35}{Problema menor do mesmo tipo}
\begin{flushright}
  \footnotesize
  \textit{\textbf{9: Funções recursivas}}\\
  \textit{c09-fundamentos-recursao-estrutural-01}\\

\end{flushright}
Em uma definição recursiva de \haskellinline|fatorial|, a chamada recursiva principal é
\haskellinline|fatorial (n - 1)| e o caso base é \haskellinline|n == 0|. Qual é a melhor interpretação dessa chamada?

\begin{enumerate}
  \item Ela é chamada apenas depois que o caso base já foi alcançado.

  \item Ela muda o problema para outro tipo de cálculo, por isso deixa de ser uma definição recursiva.

  \item Ela impede a terminação, pois toda chamada recursiva necessariamente se afasta do caso base.

  \item Ela é equivalente a uma estrutura \haskellinline|while| imperativa escrita diretamente em Haskell.

  \item Ela preserva o tipo de problema, mas o torna menor porque subtrai uma unidade de n, aproximando a computação do caso base.


\end{enumerate}

\questionTitle{36}{A variável automática it no GHCi}
\begin{flushright}
  \footnotesize
  \textit{\textbf{3: Expressões e definições}}\\
  \textit{c03-ghci-definicoes-01}\\

\end{flushright}
Considere a seguinte sessão no GHCi.

\begin{haskellcode}
> 2 + 3 * 4
14
> 7 * (it - 4)
70
> it
\end{haskellcode}

Qual valor será exibido na última linha?

\begin{enumerate}
  \item \haskellinline|2 + 3 * 4|.

  \item \haskellinline|14|.

  \item \haskellinline|70|.

  \item \haskellinline|10|.

  \item O GHCi produz erro, pois \haskellinline|it| só pode ser usado uma vez.


\end{enumerate}

\questionTitle{37}{Diagnóstico de ordem aparente}
\begin{flushright}
  \footnotesize
  \textit{\textbf{8: Programas interativos}}\\
  \textit{c08-buffering-05}\\

\end{flushright}
Um aluno observa que o terminal parece aguardar a entrada antes de mostrar o \emph{prompt}. Qual diagnóstico é mais apropriado para o seguinte padrão?

\begin{minted}{haskell}
putStr "Digite o valor: "
valor <- getLine
\end{minted}

\begin{enumerate}
  \item A ação \haskellinline|putStr| não pode ser usada antes de \haskellinline|getLine|.

  \item Haskell executou \haskellinline|getLine| antes de \haskellinline|putStr| por causa da avaliação preguiçosa.

  \item A ordem das ações está correta; o problema provável é a saída bufferizada, que pode atrasar a exibição do texto produzido por \haskellinline|putStr|.

  \item O programa está correto apenas se \haskellinline|valor| tiver tipo \haskellinline|IO String|.

  \item A variável \haskellinline|valor| deveria ser definida com \haskellinline|let|.


\end{enumerate}

\questionTitle{38}{Return não encerra um bloco do}
\begin{flushright}
  \footnotesize
  \textit{\textbf{10: Ações de E/S recursivas}}\\
  \textit{c10-return-tipos-05}\\

\end{flushright}
Muitos alunos com experiência em linguagens imperativas associam \haskellinline|return| ao encerramento imediato de uma função. Considere:

\begin{haskellcode}
interacao :: IO ()
interacao = do
  putStrLn "Início"
  return "valor intermediário"
  putStrLn "Fim"
\end{haskellcode}

Qual alternativa descreve corretamente o comportamento desse bloco?

\begin{enumerate}
  \item O programa imprime \haskellinline|"valor intermediário"| automaticamente entre as duas mensagens.

  \item O programa não compila, pois \haskellinline|return| só pode aparecer na última linha de um bloco \haskellinline|do|.

  \item A ação \haskellinline|return| extrai uma string de dentro de \haskellinline|IO|, sendo o inverso de \haskellinline|<-|.

  \item A mensagem \haskellinline|"Início"| é impressa, a ação \haskellinline|return "valor intermediário"| produz uma string que é descartada, e a mensagem \haskellinline|"Fim"| também é impressa.

  \item A execução para em \haskellinline|return "valor intermediário"|, portanto \haskellinline|"Fim"| não é impresso.


\end{enumerate}

\questionTitle{39}{Reuso uniforme e comportamento por instância}
\begin{flushright}
  \footnotesize
  \textit{\textbf{7: Polimorfismo}}\\
  \textit{cap07-comparacao-reuso-e-instancias}\\

\end{flushright}
Considere as funções:

\begin{minted}{haskell}
length :: [a] -> Int
show   :: Show a => a -> String
\end{minted}

Qual comparação é correta?

\begin{enumerate}
  \item As duas são monomórficas, pois retornam tipos concretos.
  \item \haskellinline{length} é uniformemente aplicável a listas de qualquer tipo de elemento; \haskellinline{show} depende de uma instância \haskellinline{Show} para o tipo do valor.
  \item As duas são paramétricas puras, pois ambas usam a variável de tipo \haskellinline{a}.
  \item \haskellinline{length} depende de \haskellinline{Show a}, pois precisa imprimir cada elemento antes de contar.
  \item \haskellinline{show} não usa polimorfismo, pois sempre retorna \haskellinline{String}.

\end{enumerate}

\questionTitle{40}{Definição de tipo em Haskell}
\begin{flushright}
  \footnotesize
  \textit{\textbf{4: Tipos de Dados}}\\
  \textit{c04-conceito-definicao-01}\\

\end{flushright}
Em Haskell, um \textbf{tipo} pode ser melhor definido como:

\begin{enumerate}
  \item Um comentário especial que descreve o propósito de uma função.

  \item Um nome dado a uma variável que armazena um valor numérico.

  \item Uma palavra-chave reservada para declaração de constantes.

  \item Uma função que recebe argumentos e retorna um resultado.

  \item Uma coleção de valores que compartilham propriedades e operações comuns.


\end{enumerate}

\questionTitle{41}{Mecânica da execução: Atribuição vs. Avaliação}
\begin{flushright}
  \footnotesize
  \textit{\textbf{1: Paradigmas de programação}}\\
  \textit{c01-mecanismo-reducao-01}\\

\end{flushright}
No paradigma imperativo (como na linguagem C), a computação avança predominantemente através da \textbf{execução de comandos} que alteram o estado da memória. Em contraste, qual é o mecanismo nuclear que impulsiona a computação em um programa funcional puro como Haskell?

\begin{enumerate}
  \item A modificação sequencial de valores armazenados em registradores de variáveis estáticas no escopo global do módulo compilado.

  \item A avaliação e redução de expressões através da aplicação de funções, substituindo parâmetros até que se atinja um valor normal (final).

  \item A validação de axiomas através de um motor de inferência lógica operando sobre um banco de dados de regras.

  \item A mutação explícita de estruturas de dados iteráveis através da alocação manual de ponteiros em tempo de execução.

  \item O despacho dinâmico de métodos orientados a eventos, desencadeados por interrupções do sistema operacional.


\end{enumerate}

\questionTitle{42}{Divisão inteira por subtrações sucessivas}
\begin{flushright}
  \footnotesize
  \textit{\textbf{9: Funções recursivas}}\\
  \textit{c09-projeto-divisao-inteira-01}\\

\end{flushright}
Para \haskellinline|p >= 0| e \haskellinline|q > 0|, qual ideia recursiva calcula corretamente o quociente da divisão inteira de \haskellinline|p| por \haskellinline|q|?

\begin{enumerate}
  \item Retornar sempre \haskellinline|p| quando \haskellinline|p < q|.

  \item Dividir \haskellinline|q| por \haskellinline|p| em cada chamada recursiva.

  \item Enquanto \haskellinline|p >= q|, somar \haskellinline|q| a \haskellinline|p| e subtrair \haskellinline|1| do quociente.

  \item Enquanto \haskellinline|p >= q|, subtrair \haskellinline|q| de \haskellinline|p| e somar \haskellinline|1| ao quociente.

  \item Usar o caso base \haskellinline|q == 0| para retornar o quociente.


\end{enumerate}

\questionTitle{43}{Soma com sucessor e antecessor}
\begin{flushright}
  \footnotesize
  \textit{\textbf{9: Funções recursivas}}\\
  \textit{c09-rastreamento-soma-01}\\

\end{flushright}
Considere:
\begin{haskellcode}
soma :: Integer -> Integer -> Integer
soma m n
  | n == 0    = m
  | otherwise = soma (succ m) (pred n)
\end{haskellcode}
Qual é o valor de \haskellinline|soma 4 3|?

\begin{enumerate}
  \item \haskellinline|8|

  \item \haskellinline|7|

  \item \haskellinline|1|

  \item \haskellinline|12|

  \item \haskellinline|6|


\end{enumerate}

\questionTitle{44}{Prompt com flush em programa completo}
\begin{flushright}
  \footnotesize
  \textit{\textbf{8: Programas interativos}}\\
  \textit{c08-programas-completos-05}\\

\end{flushright}
Considere o trecho:

\begin{minted}{haskell}
main :: IO ()
main = do
  putStr "Palavra: "
  hFlush stdout
  palavra <- getLine
  putStrLn palavra
\end{minted}

Qual é a finalidade de \haskellinline|hFlush stdout| nesse programa?

\begin{enumerate}
  \item Converter \haskellinline|palavra| de \haskellinline|IO String| para \haskellinline|String|.

  \item Limpar a entrada digitada pelo usuário antes de \haskellinline|getLine|.

  \item Adicionar uma quebra de linha ao final do \emph{prompt}.

  \item Imprimir automaticamente o conteúdo de \haskellinline|palavra|.

  \item Garantir que o \emph{prompt} produzido por \haskellinline|putStr| seja exibido antes de o programa aguardar a entrada com \haskellinline|getLine|.


\end{enumerate}

\questionTitle{45}{Avaliação parametrizada com precedência e associatividade}
\begin{flushright}
  \footnotesize
  \textit{\textbf{3: Expressões e definições}}\\
  \textit{c03-precedencia-associatividade-01}\\

\end{flushright}
Qual é o resultado da seguinte expressão em Haskell?

\begin{haskellcode}
10 - 5 - 2
\end{haskellcode}

\begin{enumerate}
  \item \haskellinline|10|.

  \item \haskellinline|7|.

  \item \haskellinline|13|.

  \item \haskellinline|3|.

  \item A expressão é sintaticamente inválida.


\end{enumerate}

\questionTitle{46}{Recursão para percorrer uma lista com efeitos}
\begin{flushright}
  \footnotesize
  \textit{\textbf{10: Ações de E/S recursivas}}\\
  \textit{c10-lacos-io-05}\\

\end{flushright}
Considere novamente:

\begin{haskellcode}
mostraLista :: Show a => [a] -> IO ()
mostraLista lista
  | null lista = return ()
  | otherwise = do
      print (head lista)
      mostraLista (tail lista)
\end{haskellcode}

Qual alternativa explica corretamente a repetição realizada por essa ação?

\begin{enumerate}
  \item A repetição ocorre porque \haskellinline|print| chama automaticamente \haskellinline|mostraLista| para cada elemento da lista.

  \item A repetição ocorre porque \haskellinline|null lista| percorre a lista inteira e imprime cada elemento.

  \item A função só imprime o primeiro elemento, pois a chamada recursiva é descartada por ter tipo \haskellinline|IO ()|.

  \item A lista é percorrida por \haskellinline|map| implicitamente, pois \haskellinline|IO| é uma mônada.

  \item Em cada passo, a ação imprime o primeiro elemento e chama recursivamente \haskellinline|mostraLista| com a cauda da lista.


\end{enumerate}

\questionTitle{47}{Comentários de linha e de bloco}
\begin{flushright}
  \footnotesize
  \textit{\textbf{3: Expressões e definições}}\\
  \textit{c03-comentarios-valores-01}\\

\end{flushright}
Assinale a alternativa que descreve corretamente a sintaxe de comentários em Haskell.

\begin{enumerate}
  \item Comentários são avaliados pelo GHCi como expressões do tipo \haskellinline|String|.

  \item Um comentário iniciado por \haskellinline|--| vai até o fim da linha, enquanto um comentário delimitado por \haskellinline|{-| e \haskellinline|-\}| pode ocupar várias linhas e pode ser aninhado.

  \item Comentários de bloco em Haskell não podem conter outros comentários de bloco em seu interior.

  \item Um comentário iniciado por \haskellinline|--| deve ser fechado com \haskellinline|--| antes do fim da linha.

  \item Comentários de linha e comentários de bloco só podem aparecer antes da declaração de módulo.


\end{enumerate}

\questionTitle{48}{Leitura de condicional aninhada}
\begin{flushright}
  \footnotesize
  \textit{\textbf{5: Expressão condicional}}\\
  \textit{c05-if-aninhamento-01}\\

\end{flushright}
Considere a definição:
\begin{haskellcode}
sinal n = if n < 0 then -1 else if n > 0 then 1 else 0
\end{haskellcode}
Qual é o resultado de \haskellinline|sinal 0|?

\begin{enumerate}
  \item \haskellinline|1|

  \item \haskellinline|0|

  \item \haskellinline|False|

  \item Ocorre um erro, pois o segundo \haskellinline|if| não tem \haskellinline|else|.

  \item \haskellinline|-1|


\end{enumerate}

\questionTitle{49}{Avaliação parametrizada de expressão let}
\begin{flushright}
  \footnotesize
  \textit{\textbf{3: Expressões e definições}}\\
  \textit{c03-let-escopo-01}\\

\end{flushright}
Qual é o valor da expressão abaixo?

\begin{haskellcode}
let a = 10; b = a - 6 in a + b * 2
\end{haskellcode}

\begin{enumerate}
  \item \haskellinline|18|.

  \item \haskellinline|2|.

  \item \haskellinline|14|.

  \item \haskellinline|28|.

  \item A expressão é inválida porque \haskellinline|let| só pode aparecer no início de uma definição de função.


\end{enumerate}

\questionTitle{50}{Status calculado localmente}
\begin{flushright}
  \footnotesize
  \textit{\textbf{5: Expressão condicional}}\\
  \textit{c05-where-avaliacao-01}\\

\end{flushright}
Analise a função:
\begin{haskellcode}
processa x
  | status == "par"   = x * 2
  | status == "impar" = x + 1
  where
    status = if even x then "par" else "impar"
\end{haskellcode}
Qual é o resultado de \haskellinline|processa 7|?

\begin{enumerate}
  \item Ocorre um erro, pois \haskellinline|status| só estaria visível depois da cláusula \haskellinline|where|.

  \item \haskellinline|8|

  \item Ocorre um erro de tipo na comparação com \haskellinline|"impar"|.

  \item \haskellinline|14|

  \item \haskellinline|7|


\end{enumerate}

\questionTitle{51}{Tipagem estática vs. dinâmica}
\begin{flushright}
  \footnotesize
  \textit{\textbf{4: Tipos de Dados}}\\
  \textit{c04-checagem-estatica-dinamica-01}\\

\end{flushright}
Qual é a principal diferença entre um sistema de tipos estático (como Haskell) e um dinâmico (como Python)?

\begin{enumerate}
  \item Na tipagem estática, os tipos podem mudar durante a execução; na dinâmica, são fixos.

  \item A tipagem estática é mais lenta, mas mais flexível que a dinâmica.

  \item Na tipagem estática, a verificação de tipos ocorre em tempo de compilação; na dinâmica, em tempo de execução.

  \item Não há diferença fundamental; são apenas termos diferentes para o mesmo conceito.

  \item Apenas linguagens compiladas podem ter tipagem estática.


\end{enumerate}

\questionTitle{52}{Modelagem de biblioteca}
\begin{flushright}
  \footnotesize
  \textit{\textbf{6: Estruturas de dados básicas}}\\
  \textit{c06-modelagem-01}\\

\end{flushright}
Considere as definições abaixo:
\begin{minted}{haskell}
type Titulo = String
type Autor  = String
type Ano    = Maybe Int
type Livro  = (Titulo, Autor, Ano)
\end{minted}
Qual definição representa corretamente uma biblioteca como coleção de livros?

\begin{enumerate}
  \item \haskellinline{type Biblioteca = Livro}

  \item \haskellinline{type Biblioteca = [Titulo]}

  \item \haskellinline{type Biblioteca = (Livro, Livro)}

  \item \haskellinline{type Biblioteca = Maybe Livro}

  \item \haskellinline{type Biblioteca = [Livro]}


\end{enumerate}

\questionTitle{53}{GHCup e gerenciamento de ambiente}
\begin{flushright}
  \footnotesize
  \textit{\textbf{2: Ambiente de desenvolvimento}}\\
  \textit{c02-ecossistema-ferramentas-01}\\

\end{flushright}
No ecossistema moderno da linguagem Haskell, qual é o papel exato da ferramenta \textbf{GHCup}?

\begin{enumerate}
  \item Converter automaticamente o código-fonte escrito em arquivos \texttt{.hs} para binários otimizados da linguagem C.

  \item Hospedar repositórios de código-fonte na nuvem, similar ao GitHub, mas com foco exclusivo em bibliotecas funcionais.

  \item Servir como o interpretador \emph{REPL} principal para testar e avaliar expressões matemáticas interativamente.

  \item Fornecer a interface gráfica e o editor de texto oficial recomendado pela comunidade para programação Haskell.

  \item Atuar como um instalador unificado que gerencia as versões do compilador GHC, da ferramenta de \emph{build} Cabal e do servidor de linguagem HLS.


\end{enumerate}

\questionTitle{54}{Diferença entre Float e Double}
\begin{flushright}
  \footnotesize
  \textit{\textbf{4: Tipos de Dados}}\\
  \textit{c04-tipos-float-double-01}\\

\end{flushright}
Qual é a principal diferença entre os tipos \haskellinline{Float} e \haskellinline{Double} em Haskell?

\begin{enumerate}
  \item \haskellinline{Double} tem precisão dupla, enquanto \haskellinline{Float} tem precisão simples.

  \item \haskellinline{Float} é mais rápido que \haskellinline{Double} para todos os cálculos.

  \item Não há diferença prática; são sinônimos.

  \item \haskellinline{Double} é um tipo inteiro de precisão arbitrária.

  \item \haskellinline{Float} só pode representar números inteiros.


\end{enumerate}

\questionTitle{55}{Declaração de módulo em arquivo fonte}
\begin{flushright}
  \footnotesize
  \textit{\textbf{3: Expressões e definições}}\\
  \textit{c03-arquivos-identificadores-01}\\

\end{flushright}
Qual alternativa apresenta uma declaração de módulo bem formada para o início de um arquivo Haskell?

\begin{enumerate}
  \item \haskellinline|where module Geometria|.

  \item \haskellinline|module Geometria where|.

  \item \haskellinline|import module Geometria where|.

  \item \haskellinline|module geometria where|.

  \item \haskellinline|module Geometria =|.


\end{enumerate}

\questionTitle{56}{Aplicações de missão crítica (Legado Refatorado)}
\begin{flushright}
  \footnotesize
  \textit{\textbf{1: Paradigmas de programação}}\\
  \textit{c01-engenharia-software-01}\\

\end{flushright}
Você foi encarregado de arquitetar o \emph{backend} de um sistema bancário de larga escala. O sistema processa milhões de transações diárias que passam por validações, conversões e cálculos de risco. A \textbf{correção matemática} e a \textbf{auditabilidade} do fluxo de dados são características críticas inegociáveis.

Entre as opções, qual paradigma de programação fornece as abstrações arquiteturais mais seguras para este domínio específico e por quê?

\begin{enumerate}
  \item O paradigma funcional, porque a imutabilidade estrita e a ausência de efeitos colaterais tornam o fluxo de dados totalmente explícito e previsível, facilitando a auditoria e as provas lógicas.

  \item O paradigma puramente orientado a objetos, porque ocultar transações dentro de objetos de estado opaco garante que auditores externos não consigam corromper os dados acidentalmente.

  \item O paradigma imperativo procedural, porque a manipulação direta e mutável do estado da memória é a única forma técnica de garantir a performance de I/O exigida por bancos de dados modernos.

  \item A escolha do paradigma é completamente irrelevante para características de software como auditabilidade e correção matemática; essas métricas dependem exclusivamente da linguagem SQL utilizada.

  \item O paradigma lógico, porque todo o ecossistema bancário pode ser facilmente abstraído como um banco de dados Prolog e resolvido instantaneamente via inferência de axiomas.


\end{enumerate}

\questionTitle{57}{Composição de take, drop e sum}
\begin{flushright}
  \footnotesize
  \textit{\textbf{6: Estruturas de dados básicas}}\\
  \textit{c06-listas-operacoes-01}\\

\end{flushright}
Analise a expressão abaixo.
\begin{minted}{haskell}
let xs = [1..12]
    ys = take 4 (drop 3 xs)
in sum ys
\end{minted}
Qual é o valor final calculado?

\begin{enumerate}
  \item A expressão resulta em erro.

  \item \haskellinline{ 30 }

  \item \haskellinline{ 22 }

  \item \haskellinline{ 34 }

  \item \haskellinline{ 78 }


\end{enumerate}

\questionTitle{58}{Avaliação básica de guardas}
\begin{flushright}
  \footnotesize
  \textit{\textbf{5: Expressão condicional}}\\
  \textit{c05-guardas-avaliacao-01}\\

\end{flushright}
Considere a função:
\begin{haskellcode}
f x
  | x < 0 = -1
  | x == 0 = 0
  | otherwise = 1
\end{haskellcode}
Qual é o valor de \haskellinline|f 0|?

\begin{enumerate}
  \item \haskellinline|0|

  \item \haskellinline|-1|

  \item \haskellinline|1|

  \item \haskellinline|"0"|

  \item Ocorre um erro de guardas não exaustivas.


\end{enumerate}

\questionTitle{59}{Quando guardas são mais idiomáticas}
\begin{flushright}
  \footnotesize
  \textit{\textbf{5: Expressão condicional}}\\
  \textit{c05-if-vs-guardas-01}\\

\end{flushright}
Em qual situação o uso de guardas tende a ser mais idiomático e legível do que \haskellinline|if-then-else| aninhado?

\begin{enumerate}
  \item Quando há apenas um teste booleano simples, com dois resultados possíveis.

  \item Quando a função não recebe argumentos.

  \item Quando os ramos do \haskellinline|if| têm tipos diferentes.

  \item Quando se deseja evitar qualquer comparação booleana.

  \item Quando a função trata vários casos sequenciais, cada um coberto por uma guarda.


\end{enumerate}

\questionTitle{60}{getLine seguido de read}
\begin{flushright}
  \footnotesize
  \textit{\textbf{8: Programas interativos}}\\
  \textit{c08-conversao-validacao-05}\\

\end{flushright}
Qual trecho expressa corretamente a ideia de ler uma linha e convertê-la para \haskellinline|Int|?

\begin{enumerate}
  \item \begin{minted}{haskell}
let linha = getLine
let n = read linha :: Int
\end{minted}

  \item \begin{minted}{haskell}
linha <- read getLine
let n = linha :: Int
\end{minted}

  \item \begin{minted}{haskell}
n <- read getLine :: IO Int
\end{minted}

  \item \begin{minted}{haskell}
let n <- getLine :: Int
\end{minted}

  \item \begin{minted}{haskell}
linha <- getLine
let n = read linha :: Int
\end{minted}


\end{enumerate}

\questionTitle{61}{Regra de layout em do}
\begin{flushright}
  \footnotesize
  \textit{\textbf{8: Programas interativos}}\\
  \textit{c08-blocos-do-05}\\

\end{flushright}
Em um bloco \haskellinline|do| escrito sem chaves e ponto e vírgula explícitos, qual regra de layout é essencial?

\begin{enumerate}
  \item Cada linha deve começar obrigatoriamente na primeira coluna do arquivo.

  \item As ações devem ficar cada vez mais indentadas à medida que aparecem no bloco.

  \item As ações pertencentes ao mesmo bloco devem ficar alinhadas na mesma coluna de indentação, a partir do primeiro item do bloco.

  \item A indentação é ignorada em Haskell; apenas quebras de linha importam.

  \item A primeira ação do bloco deve ficar menos indentada que a palavra \haskellinline|do|.


\end{enumerate}

\questionTitle{62}{Inferência da média de uma lista}
\begin{flushright}
  \footnotesize
  \textit{\textbf{7: Polimorfismo}}\\
  \textit{cap07-inferencia-media-lista}\\

\end{flushright}
Qual é a assinatura de tipo mais geral da função abaixo?

\begin{minted}{haskell}
media xs = sum xs / fromIntegral (length xs)
\end{minted}

\begin{enumerate}
  \item \mintinline{haskell}{media :: Integral a => [a] -> a}
  \item \mintinline{haskell}{media :: [Int] -> Int}
  \item \mintinline{haskell}{media :: RealFrac a => [a] -> Int}
  \item \mintinline{haskell}{media :: Fractional a => [a] -> a}
  \item \mintinline{haskell}{media :: Num a => [a] -> a}

\end{enumerate}

\questionTitle{63}{Restrições combinadas em função sobrecarregada}
\begin{flushright}
  \footnotesize
  \textit{\textbf{7: Polimorfismo}}\\
  \textit{cap07-adhoc-funcao-inferida-num-show}\\

\end{flushright}
Considere a definição:

\begin{minted}{haskell}
f x = show (x + 1)
\end{minted}

Qual assinatura de tipo é a mais geral?

\begin{enumerate}
  \item \mintinline{haskell}{f :: Show a => a -> String}
  \item \mintinline{haskell}{f :: (Num a, Show a) => a -> String}
  \item \mintinline{haskell}{f :: a -> String}
  \item \mintinline{haskell}{f :: Num a => a -> String}
  \item \mintinline{haskell}{f :: Int -> Int}

\end{enumerate}

\questionTitle{64}{Assinatura de função com dois Int e String}
\begin{flushright}
  \footnotesize
  \textit{\textbf{4: Tipos de Dados}}\\
  \textit{c04-funcao-assinatura-01}\\

\end{flushright}
Qual das seguintes assinaturas de tipo representa uma função que recebe dois \haskellinline{Int} e uma \haskellinline{String} e retorna um \haskellinline{Bool}?

\begin{enumerate}
  \item \haskellinline{(Int, Int, String) -> Bool}

  \item \haskellinline{Int -> Int -> String, Bool}

  \item \haskellinline{Bool -> Int -> Int -> String}

  \item \haskellinline{[Int, Int, String] -> Bool}

  \item \haskellinline{Int -> Int -> String -> Bool}


\end{enumerate}

\questionTitle{65}{Enum e Bounded não são equivalentes}
\begin{flushright}
  \footnotesize
  \textit{\textbf{7: Polimorfismo}}\\
  \textit{cap07-hierarquia-enum-bounded}\\

\end{flushright}
Qual afirmação distingue corretamente as classes \haskellinline{Enum} e \haskellinline{Bounded}?
\begin{enumerate}
  \item Todo tipo instância de \haskellinline{Enum} é necessariamente instância de \haskellinline{Bounded}.
  \item As duas classes pertencem apenas à hierarquia de \haskellinline{Fractional}.
  \item \haskellinline{Enum} descreve tipos que podem ser enumerados com operações como \haskellinline{succ} e \haskellinline{pred}; \haskellinline{Bounded} descreve tipos com \haskellinline{minBound} e \haskellinline{maxBound}.
  \item \haskellinline{Enum} e \haskellinline{Bounded} são a mesma classe com nomes diferentes.
  \item \haskellinline{Bounded} permite usar \haskellinline{succ}, enquanto \haskellinline{Enum} fornece \haskellinline{minBound}.

\end{enumerate}

\questionTitle{66}{Valor opcional}
\begin{flushright}
  \footnotesize
  \textit{\textbf{6: Estruturas de dados básicas}}\\
  \textit{c06-maybe-conceito-01}\\

\end{flushright}
Qual opção representa corretamente a ideia central da estrutura \haskellinline{Maybe}?

\begin{enumerate}
  \item Ela representa listas de tamanho variável.

  \item Ela representa explicitamente a presença ou a ausência de um valor.

  \item Ela representa tuplas com campos mutáveis.

  \item Ela substitui o tipo \haskellinline{Bool} em expressões condicionais.

  \item Ela é usada apenas para números opcionais.


\end{enumerate}

\questionTitle{67}{Refatoração para acumulador}
\begin{flushright}
  \footnotesize
  \textit{\textbf{9: Funções recursivas}}\\
  \textit{c09-cauda-refatoracao-01}\\

\end{flushright}
Qual alternativa é uma refatoração correta de \haskellinline|fatorial| usando uma função auxiliar com acumulador?

\begin{enumerate}
  \item \begin{haskellcode}
fat n = fat' n 0
  where
    fat' n acc
      | n == 0 = acc
      | n > 0  = fat' (n - 1) (n * acc)
\end{haskellcode}

  \item \begin{haskellcode}
fat n = fat' n 1
  where
    fat' n acc
      | n == 0 = acc
      | n > 0  = n * fat' (n - 1) acc
\end{haskellcode}

  \item \begin{haskellcode}
fat n = fat' n 1
  where
    fat' n acc
      | n == 0 = 1
      | n > 0  = fat' (n + 1) (n * acc)
\end{haskellcode}

  \item \begin{haskellcode}
fat n = fat' n 1
  where
    fat' n acc
      | n == 0 = n
      | n > 0  = fat' (n - 1) acc
\end{haskellcode}

  \item \begin{haskellcode}
fat n = fat' n 1
  where
    fat' n acc
      | n == 0 = acc
      | n > 0  = fat' (n - 1) (n * acc)
\end{haskellcode}


\end{enumerate}

\questionTitle{68}{Consequência da falta de cobertura}
\begin{flushright}
  \footnotesize
  \textit{\textbf{5: Expressão condicional}}\\
  \textit{c05-guardas-exaustividade-01}\\

\end{flushright}
O que acontece se uma função com guardas é chamada com uma entrada para a qual nenhuma guarda resulta em \haskellinline|True|, e não existe outra equação aplicável?

\begin{enumerate}
  \item O compilador sempre rejeita o programa com erro de compilação.

  \item A função passa a retornar automaticamente \haskellinline|undefined|.

  \item O avaliador repete a verificação das guardas até alguma se tornar verdadeira.

  \item Ocorre um erro em tempo de execução, porque a definição não cobre aquela entrada.

  \item A função retorna o valor da última expressão escrita na definição.


\end{enumerate}

\questionTitle{69}{Paradigmas - Característica Fundamental (Legado Refatorado)}
\begin{flushright}
  \footnotesize
  \textit{\textbf{1: Paradigmas de programação}}\\
  \textit{c01-conceitos-paradigmas-01}\\

\end{flushright}
Qual das seguintes características é a mais fundamental para definir uma linguagem como \textbf{puramente funcional}?

\begin{enumerate}
  \item A ausência completa de um sistema de tipos estático.

  \item A obrigatoriedade de usar o paradigma de orientação a objetos para gerenciar o estado da aplicação.

  \item O tratamento de funções como valores de primeira classe e a total ausência de efeitos colaterais.

  \item A capacidade de usar laços imperativos nativos como \cinline|for| e \cinline|while|.

  \item O uso de ponteiros abstratos para a manipulação direta e manual da memória pelo programador.


\end{enumerate}


\end{multicols}
\makeanswersheet{UFOP}{BCC222: Programação Funcional}{José Romildo}{2026/1}{Primeira Prova}{01}{69}{25}{7}
\gabarito{01}{C,B,C,A,B,B,B,E,A,C,C,C,D,D,E,D,E,C,D,B,C,C,B,E,E,D,B,A,E,B,D,C,C,C,E,C,C,D,B,E,B,D,B,E,D,E,B,B,A,B,C,E,E,A,B,A,B,A,E,E,C,D,B,E,C,B,E,D,C}
\end{document}
